{\large\textbf{Waves and Phase Transitions in Spin Systems (10.0 points)}}

\textbf{Introduction}

In classical physics, angular momentum arises from the motion of an object
around an axis - whether it be a spinning top, a rotating planet, or an
orbiting electron in the atom. However, in quantum physics, fundamental
particles possess an intrinsic and quantized form of angular momentum called
\textit{spin}. This property plays a crucial role in various physical
phenomena, ranging from materials properties, such as magnetism, to modern
applications, such as quantum computing.

In this problem we will treat spin classically, which will lead to some
qualitatively correct results. You will explore the physics of spin systems
through spin-spin interactions, evolution under magnetic fields, and
statistical physics to understand the emergence of spin waves and phase
transitions in magnets.

\textbf{Useful information:}

$\cosh(x) \equiv \frac{e^x + e^{-x}}{2}$, $\sinh(x) \equiv \frac{e^x - e^{-x}}{2}$,
$\tanh(x) \equiv \frac{\sinh(x)}{\cosh(x)} \approx x - \frac{1}{3}x^3$ for
$|x| \ll 1$

The magnetic field due to a magnetic dipole of moment $\vec{\mu}$ at a position
$\vec{r}$ away from it is given by ($\mu_0$ is is the vacuum permeability):
\begin{equation}
\vec{B} = \frac{\mu_o}{4\pi} \left( \frac{3(\vec{\mu} \cdot \vec{r})\vec{r}}{r^5} - \frac{\vec{\mu}}{r^3} \right) \tag{1}
\end{equation}

\textbf{Part A. Precession and interactions of magnetic dipoles (1.2 points)}

Consider a ring of radius $R$, total mass $M$, and charge $Q > 0$ distributed
uniformly. The ring rotates with an angular speed $\omega$ around a
perpendicular axis that passes through its center of mass.

\textbf{A.1}\quad It is possible to write the ring's magnetic moment
$\vec{\mu}$ in terms of its angular momentum $\vec{L}$ as
$\vec{\mu} = \gamma\vec{L}$. Find the constant $\gamma$, called the
\textit{gyromagnetic ratio}, of this system in terms of $Q$ and $M$.
\hfill 0.3 pt

The ring is placed in a weak uniform magnetic field $\vec{B} = B\hat{z}$ ,
making an angle $\theta$ with $\vec{\omega}$, see \textit{Figure A.1}.

\begin{center}
\includegraphics[width=0.29\linewidth]{figures/APhO_2025_Q2_fig1.png}\\
\textit{Figure A.1.}
\end{center}

\textbf{A.2}\quad Find the angular frequency $\omega_L$ of the angular momentum
precession (the so-called Larmor frequency) due to the external magnetic field
in terms of $B$ and $\gamma$. Take the positive direction to be
counter-clockwise with respect to $+z$. \hfill 0.4 pt

Now we turn off the external magnetic field and place an identical ring at a
horizontal distance $d \gg R$ from the original ring such that the magnetic
moment of the new ring $\vec{\mu}_2$ makes an angle $\theta$ with
$\vec{\mu}_1$, see \textit{Figure A.2}.

\begin{center}
\includegraphics[width=0.52\linewidth]{figures/APhO_2025_Q2_fig2.png}\\
\textit{Figure A.2.}
\end{center}

\textbf{A.3}\quad The magnetic interaction energy between the two rings can be
written as $U = J_0\vec{L}_1 \cdot \vec{L}_2$, where $J_0$ is a constant and
$\vec{L}_i$ is the angular momentum of the $i$th ring. Find $J_0$ in terms of
$\gamma, d$ and fundamental constants. \hfill 0.5 pt

\textbf{Part B. Spin Waves (4.5 points)}

In what follows we investigate the dynamics of spins. A spin is a particle with
intrinsic angular momentum $\vec{S}$, which has an associated magnetic moment
$\vec{\mu}$ related to $\vec{S}$ via the gyromagnetic ratio as in
\textbf{Part A.1}, $\vec{\mu} = \gamma\vec{S}$.

The magnetic dipoles of two spins interact with each other. However, this
interaction is negligible compared to another interaction arising from a
quantum mechanical origin, which is not present in classical systems.
Interestingly, the energy associated with this quantum interaction has the same
form which we found in \textbf{Part A.3}, scaling with
$\vec{S}_1 \cdot \vec{S}_2$, albeit with the opposite sign.

Now we will look at a very long chain of spins. The positions of the spins are
fixed along the $x$-axis, with a distance $a$ separating them, see
\textit{Figure B.1}. We will approximate the total energy of the system by
considering the interactions between nearest neighbors only, so that the
energy can be written as
\[
E = -J \sum_i \vec{S}_i \cdot \vec{S}_{i+1}
\]
where $J > 0$ is the interaction strength, and $\vec{S}_i$ is the spin angular
momentum vector of the $i$th dipole, with magnitude $S$. The spin vectors are
free to rotate in three dimensions. Notice that the sign of the energy is
different from the last part. This interaction is purely quantum mechanical.

\begin{center}
\includegraphics[width=0.57\linewidth]{figures/APhO_2025_Q2_fig3.png}\\
\textit{Figure B.1.}
\end{center}

\textbf{B.1}\quad The energy terms containing $\vec{S}_i$ in the sum above can
be viewed as the interaction energy between an effective magnetic field
$\vec{B}_{i,\mathrm{eff}}$ and the magnetic moment of $\vec{S}_i$. Find
$\vec{B}_{i,\mathrm{eff}}$ and express your answer in terms of $J$, the
gyromagnetic ratio $\gamma$, and other spins $\vec{S}_j$ (specify the indices
$j$ in relation to $i$) \hfill 0.3 pt

\textbf{B.2}\quad Using the concept of effective magnetic field, express the
rate of change of the $i$th spin vector, $d\vec{S}_i/dt$, in terms of
$J, \vec{S}_i$, and other spins $\vec{S}_j$ (specify the indices $j$ in
relation to $i$). \hfill 0.3 pt

For the rest of \textbf{Part B}, assume that the system is highly magnetized
along the $z$ direction, so we can use the approximations $S_{i,z} \approx S$
and $dS_{i,z}/dt \approx 0$ for each spin, see \textit{Figure B.2}. In this
regime, the set of equations describing the spins time evolution is satisfied
by a traveling wave solution for $S_{i,x}$and $S_{i,y}$ characterized by a wave
vector $k$ and angular frequency $\omega$.

\begin{center}
\includegraphics[width=0.49\linewidth]{figures/APhO_2025_Q2_fig4.png}\\
\textit{Figure B.2.}
\end{center}

\textbf{B.3}\quad Find the relationship between $\omega$ and $k$ (known as the
dispersion relation, $\omega(k)$) for the spin waves in terms of $J, S$ and
$a$. \textit{Hint}: express the position of the $i$th spin as $x = a \cdot i$.
\hfill 2.0 pt

The spin wave described above carries energy and momentum. At low energies,
the relation between its energy and momentum resembles that of a massive
classical particle with an effective mass $m_{\mathrm{eff}}$ , a concept known
as a \textit{quasi-particle}.

\textbf{B.4}\quad For small $k$ ($k \ll 1/a$), find the effective mass
$m_{\mathrm{eff}}$ of the spin wave. Express your answer in terms of $J, S, a$
and fundamental constants. \hfill 0.6 pt

Spin waves can be experimentally probed using inelastic neutron scattering.
Although neutrons have zero net charge, they have a finite spin, allowing them
to interact with other spins.

\textbf{B.5}\quad Suppose that initially, all the spins in the chain are
pointing along the $z$ direction. A neutron with low energy travels on the
$x{-}y$ plane making an incident angle $\theta_{in}$ with the chain and
scatters with an angle $\theta_{out}$ as shown in \textit{Figure B.3}.
Assuming the neutron excites a single low wave vector spin wave, find the
effective mass $m_{\mathrm{eff}}$ of the spin wave, in terms of
$\theta_{\mathrm{in}}, \theta_{\mathrm{out}}$ and the neutron mass $m_n$.
Assume that the chain stays at rest. \hfill 1.3 pt

\begin{center}
\includegraphics[width=0.49\linewidth]{figures/APhO_2025_Q2_fig5.png}\\
\textit{Figure B.3.}
\end{center}

\textbf{Part C. Phase transitions in spin chains (4.3 points)}

Next we consider the same chain made of $N$ spins from \textbf{Part B}, except
the spin vectors are now restricted to point either up or down along the
$z$-axis, so that the spin component along $z$ can be written as
$S_{i,z} = s_i S$, where $s_i = \pm 1$, see \textit{Figure C.1}. In addition to
the nearest neighbor interactions, we could have an external magnetic field
pointing along the $z$-axis so that the total energy of the system is given by
\[
E = -\tilde{J} \sum_i s_i s_{i+1} - h \sum_i s_i .
\]
We assume $\tilde{J} \geq 0$, and $h$ is a constant dependent on the magnetic
field. The spin system is at equilibrium with a heat bath at temperature $T$.
Ignore the edges of the chain.

\begin{center}
\includegraphics[width=0.51\linewidth]{figures/APhO_2025_Q2_fig6.png}\\
\textit{Figure C.1.}
\end{center}

\textbf{C.1}\quad Assume first that $\tilde{J} = 0$, what is the ratio between
the probability to find an arbitrary spin aligned to the magnetic field
$p_\uparrow$ to being anti-aligned to the magnetic field $p_\downarrow$?
Express $p_\uparrow/p_\downarrow$ in terms of $h$, $T$ and fundamental
constants. \hfill 0.5 pt

\textbf{C.2}\quad Find the average polarization of the system
$\bar{s} \equiv \frac{1}{N}\sum_i s_i$ for $N \gg 1$ in terms of $h$, $T$ and
fundamental constants. If the magnetic field $h$ can range from $-h_0$ to
$h_0$, make a sketch of $\bar{s}$ as a function of $h$ for the cases
$h_o \gg k_B T$, $h_o \approx k_B T$ and $h_o \ll k_B T$. \hfill 1.0 pt

In the remaining questions, we turn off the magnetic field, so $h = 0$, and set
$\tilde{J} > 0$.

\textbf{C.3}\quad What is the energy $E_g$ of the ground state (the lowest
energy state)? Express your answer in terms of $\tilde{J}$ and $N$.
\hfill 0.2 pt

Instead of considering the interactions between each spin and its neighbors,
we assume that each spin sees an average polarization $\bar{s}$ from its
nearest-neighbors.

\textbf{C.4}\quad Approximate the energy of the system as a sum over all spins
\hfill 0.2 pt
\[
E = -\tilde{J}_{\mathrm{eff}} \sum_i s_i
\]
and express $\tilde{J}_{\mathrm{eff}}$in terms of $\tilde{J}$ and $\bar{s}$.

\textbf{C.5}\quad Using your result from \textbf{C.2}, find an equation that the
average polarization $\bar{s}$ must satisfy. The number of solutions to this
equation depends on $T$. Find the critical temperature $T_c$ at which the
number of solutions changes. Express your answer in terms of $\tilde{J}$ and
fundamental constants. \hfill 1.2 pt

\textbf{C.6}\quad Find all possible values of $\bar{s}$ when $T < T_c$ and
$T_c - T \ll T_c$. Express your answers in terms of $T$ and $T_c$. Sketch all
possible values of $\bar{s}$ for the temperature $T$ in the range
$0 \leq T \leq 2T_c$. \hfill 1.0 pt

\textbf{C.7}\quad What magnetic phase of matter does $T > T_c$correspond to?
How about when $T < T_c$? Choose between paramagnetic or ferromagnetic.
\hfill 0.2 pt
